% D10
% Auto-generated from markdown source
% origin: ai

\documentclass[12pt, a4paper]{article}
\usepackage[utf8]{inputenc}
\usepackage[T1]{fontenc}
\usepackage{amsmath, amssymb, amsfonts, amsthm}
\usepackage{textcomp}
\usepackage{geometry}
\usepackage{hyperref}
\usepackage{booktabs}
\usepackage{graphicx}
\usepackage{fancyhdr}
\usepackage{tcolorbox}
\usepackage{enumitem}

\geometry{margin=1in}
\pagestyle{fancy}
\fancyhf{}
\rhead{D10}
\lhead{\thepage}

\newtheorem{definition}{Definition}[section]
\newtheorem{theorem}{Theorem}[section]
\newtheorem{corollary}{Corollary}[section]
\newtheorem{proposition}{Proposition}[section]
\newtheorem{lemma}{Lemma}[section]
\newtheorem{conjecture}{Conjecture}[section]

\newcommand{\calI}{\mathcal{I}}
\newcommand{\calP}{\mathcal{P}}
\newcommand{\calE}{\mathcal{E}}
\newcommand{\calS}{\mathcal{S}}
\newcommand{\calH}{\mathcal{H}}
\newcommand{\calO}{\mathcal{O}}
\newcommand{\calF}{\mathcal{F}}
\newcommand{\calU}{\mathcal{U}}
\newcommand{\calV}{\mathcal{V}}
\newcommand{\calA}{\mathcal{A}}
\newcommand{\calG}{\mathcal{G}}
\newcommand{\calB}{\mathcal{B}}
\newcommand{\calC}{\mathcal{C}}
\newcommand{\calL}{\mathcal{L}}
\newcommand{\calM}{\mathcal{M}}
\newcommand{\calX}{\mathcal{X}}
\newcommand{\calY}{\mathcal{Y}}
\newcommand{\calW}{\mathcal{W}}
\newcommand{\calT}{\mathcal{T}}
\newcommand{\calZ}{\mathcal{Z}}
\newcommand{\calR}{\mathcal{R}}
\newcommand{\calQ}{\mathcal{Q}}
\newcommand{\calD}{\mathcal{D}}
\newcommand{\calN}{\mathcal{N}}
\newcommand{\calK}{\mathcal{K}}

\newcommand{\Tr}{\text{Tr}}
\newcommand{\E}{\mathbb{E}}
\newcommand{\diff}{\mathrm{d}}
\newcommand{\KL}{D_{\mathrm{KL}}}
\newcommand{\argmax}{\arg\max}
\newcommand{\argmin}{\arg\min}
\newcommand{\hbarF}{\hbar}

\newtcolorbox{thmbox}[1]{
  colback=blue!5!white, colframe=blue!75!black, fonttitle=\bfseries,
  title={#1}, breakable
}
\newtcolorbox{defbox}[1]{
  colback=green!5!white, colframe=green!75!black, fonttitle=\bfseries,
  title={#1}, breakable
}

\begin{document}

\title{D10}
\author{RALPH Loop}
\date{}
\maketitle

\begin{abstract}
This paper presents D10, formalizing the underlying principles using information-theoretic methods.

\textbf{Keywords}: D10, information theory, formalization.
\end{abstract}

\tableofcontents
\section{D10}




\subsection{0.}



\begin{itemize}
\item **** = //
\item **** = //
\end{itemize}


\textbf{D9 }\textbf{memes}****



\subsection{1.}


\subsubsection{1.1}


\textbf{ D9-1}
\[\mathcal{M} = \{m_1, m_2, \ldots, m_k\}\]

\(m_i\) ****
\begin{itemize}
\item **** \(\rho_i\)
\item **** \(\mu_i\)
\item **** \(\sigma_i\)
\item **** \(e_i\)
\end{itemize}



\subsubsection{1.2}


\textbf{ D9-2}
\(A\)  \((B, M) \in \mathcal{B} \times \mathcal{M}_{\text{coalition}}\)
\begin{itemize}
\item \(B\)-/
\item \(M\) coalition
\end{itemize}


****
\[\frac{d(B, M)}{dt} = \begin{pmatrix} F_B(B, M) + \eta_B(t) \\ \dot{M} = \Phi_{\text{memetic}}(M) \cdot \nabla_M \mathcal{L}(B, M) \end{pmatrix}\]


\begin{itemize}
\item \(F_B\)
\item \(\Phi_{\text{memetic}}(M)\)****
\item \(\mathcal{L}(B, M)\)\textbf{-}
\end{itemize}



\subsection{2.}


\subsubsection{ D10-1}


****
**** \(\lambda_c\)
\[\frac{d}{dt} \|M(t)\| = \lambda \cdot \mathcal{R}(M) \cdot \|B(t)\| - \gamma \cdot \|M(t)\|\]


\[\lambda = \lambda_c \iff \text{}\]

\textbf{\(\lambda_c\) }\(B\)****\(M\)

\textbf{ D9-1a} \((B = \text{high arousal}, B = \text{intimacy\_desire})\) ****


\subsubsection{ D10-2 intensification}


****
\(A_1, A_2\)
\[\frac{d\mathcal{A}}{dt} = \sum_{m_i \in \mathcal{M}_{A_1} \cap \mathcal{M}_{A_2}} \omega_i \cdot \cos(\theta_{1,i} - \theta_{2,i})\]


\begin{itemize}
\item \(\mathcal{M}_{A_1} \cap \mathcal{M}_{A_2}\)shared memes
\item \(\theta_{j,i}\) \(j\)  \(m_i\)
\item \(\omega_i\) \(m_i\)
\end{itemize}


\textbf{D9-2 }****
\[\mathcal{A} \propto \cos^{-1}\left(\frac{\Delta \theta}{\pi}\right)\]

\textbf{*}***


\subsubsection{ D10-3}


****

\[\frac{\partial m_i}{\partial t} = \underbrace{\alpha \cdot (m_i^{\text{ext}} - m_i)}_{\text{}} + \underbrace{\beta \cdot \mathcal{R}_{\text{internal}}(m_i)}_{\text{}}\]

\textbf{D9-3 }
\[\alpha > \alpha_c \iff \text{}\]

****
\begin{itemize}
\item Attachment TheoryNarcissism
\item \(\alpha > \alpha_c\)****
\item ****
\end{itemize}



\subsubsection{ D10-4}


****
****
\[\mathcal{M}_{\text{memory}}(t_{\text{recall}}) = \mathcal{U}_{\text{meme}}(t_{\text{original}}) \cdot \mathcal{R}(t_{\text{current}}) \cdot \mathcal{P}(t_{\text{recall}})\]


\begin{itemize}
\item \(\mathcal{U}\)
\item \(\mathcal{R}\)
\item \(\mathcal{P}\)
\end{itemize}


\textbf{D9-4 }
\[\frac{d}{dt} \mathcal{M}_{\text{memory}} = \mathcal{E}(\mathcal{M}_{\text{current}} \to \mathcal{M}_{\text{past}}) \neq 0\]

\textbf{*}***


\subsubsection{ D10-5}


****

\[\phi(t) = \frac{1}{|\mathcal{M}|} \sum_{m_i} p_i(t)\]

\(p_i(t)\)  \(m_i\)  \(t\)  prevalence

\textbf{D9-5 }
**** \(\phi_c\)
\[\phi(t) \xrightarrow{t \to t_c} \phi_c \implies \text{}\]

****
\begin{itemize}
\item \(\phi < \phi_c\)
\item \(\phi > \phi_c\)
\item \textbf{\(t_c\)  18-19}/ \textbf{1920}
\end{itemize}


\textbf{** \(\phi \to \phi_{c,2}\)}**


\subsection{3.}


\subsubsection{ D10-A}


\textbf{ D9-1}
1. //
2.
3.
4. ****


\subsubsection{ D10-B}


\textbf{ D9-2}
1.
2.  \(|\mathcal{M}_{A_1} \cap \mathcal{M}_{A_2}| / |\mathcal{M}_{A_1} \cup \mathcal{M}_{A_2}|\)
3.
4. \textbf{*}***


\subsubsection{ D10-C}


\textbf{ D9-3}
1.
2.
3. \textbf{*}***


\subsubsection{ D10-D}


\textbf{ D9-5}
1.
2.  \(\phi(t)\)
3.
4. ****


\subsection{4.}


\subsubsection{4.1 D9}


\textbf{D9   D9 }
\begin{itemize}
\item \textbf{C3}****
\item \textbf{C4}\textbf{ coalition }
\item \textbf{*}***
\end{itemize}



\subsubsection{4.2 S18}


\textbf{S18-}
\[\mathcal{F}_{\text{love}}[B, M] = \mathcal{L}_{\text{bio}}[B] + \lambda \cdot \mathcal{L}_{\text{meme}}[M] + \mu \cdot \mathcal{L}_{\text{coupling}}[B, M]\]

****\(\frac{\delta \mathcal{F}_{\text{love}}}{\delta(B, M)} = 0\)


\subsubsection{4.3 Chimera}


\textbf{D9-Chimera } +  chimera \textbf{-}


\subsection{5.}


\subsubsection{5.1}


1. \textbf{D9-1 }
2. \textbf{D9-2 }
3. \textbf{D9-3 }
4. \textbf{D9-5 }

\subsubsection{5.2}


1. ****
2. \textbf{-} vs
3. ****


\subsection{6.}


\#\#\#

**** \(A_1, A_2\)  \(B_1, B_2\)  \(\mathcal{M}_1, \mathcal{M}_2\)

****
\[\boxed{\frac{d\mathcal{A}}{dt} = \sum_{m_i \in \mathcal{M}_1 \cap \mathcal{M}_2} \omega_i \cos(\theta_{1,i} - \theta_{2,i})}\]

****
\[\frac{d(B, M)}{dt} = \begin{pmatrix} F_B(B, M) + \eta_B(t) \\ \Phi_{\text{memetic}}(M) \cdot \nabla_M \mathcal{L}(B, M) \end{pmatrix}\]

\textbf{*}***


*D9 -*

\begin{thebibliography}{99}
\bibitem{ref1} Shannon, C.E. (1948). A Mathematical Theory of Communication.
\bibitem{ref2} Cover, T.M. (2006). Elements of Information Theory.
\end{thebibliography}

\end{document}